\documentclass[12pt,letterpaper]{article}

% ---- Packages ----
\usepackage[utf8]{inputenc}
\usepackage[T1]{fontenc}
\usepackage{amsmath,amssymb}
\usepackage{newtxtext,newtxmath}  % Times New Roman text + math (replaces mathptmx)
\usepackage[margin=1in]{geometry}
\usepackage{setspace}
\usepackage{titlesec}
\usepackage{enumitem}
\usepackage{xr-hyper}
\usepackage{hyperref}
\externaldocument{1 - General Overview}
\externaldocument{7 - Quantum Mechanics}
\usepackage{microtype}
\usepackage{booktabs}

% ---- Section formatting ----
\titleformat{\section}{\large\bfseries\sffamily}{\thesection.}{0.5em}{}
\titleformat{\subsection}{\normalsize\bfseries\sffamily}{\thesubsection}{0.5em}{}
\titleformat{\subsubsection}{\normalsize\bfseries}{\thesubsubsection}{0.5em}{}

% ---- Spacing ----
\setstretch{1.15}
\setlength{\parskip}{0.5em}
\setlength{\parindent}{0em}

% ---- Custom commands ----
\newcommand{\rhoa}{\rho_{\!A}}
\newcommand{\Ka}{K_{\!A}}
\newcommand{\vin}{v_{\mathrm{in}}}
\newcommand{\rs}{r_s}
% \hbar is already defined in LaTeX; no redefinition needed

\begin{document}

\begin{center}
{\LARGE\bfseries Unifying Theory of Dynamic Aether Mechanics:\\
Strong and Weak Nuclear Forces}\\[1em]
{\large Erik Otto}
\end{center}

\vspace{1em}

\begin{abstract}
Within the Dynamic Aether Mechanics framework, the strong force arises from
compressed Aether bridges---physical strands of quantum superfluid compressed to
$\sim\!7 \times 10^7$ times ambient density, connecting quark vortex endpoints.
This paper derives confinement (bridge snapping at $\sim$1.2~fm creates new
quark--antiquark pairs), asymptotic freedom (bridges are ``all core'' with
compactness $d/\xi \approx 1$--$2$), hadron mass (predominantly bridge compression
energy), and the deconfinement phase transition at $T \approx 170$~MeV.  Quantum
fluctuations of the bridge reproduce the L\"uscher term $-\pi/(12\,r)$, predict
Schwinger pair-creation rates that naturally separate chiral from potential-model
QCD at a crossover mass of $\sim$250~MeV, and explain the nucleon sigma term.
The weak force is identified as vortex topology transitions.  A unified summary
table relates all four fundamental forces through distinct Aether processes.
\end{abstract}

\vspace{1em}

\bigskip
\noindent\textit{This paper is part of a series presenting the Unifying Theory of
Dynamic Aether Mechanics.  For the foundational concepts of Aether binding,
the three independent properties (density, pressure, temperature), and the
unification of force constants, see Paper~1: General Overview \cite{otto-overview}.}
\bigskip

\section{The Strong Force: Compressed Aether Bridges}
\label{sec:strong-force}
% ============================================================

The three forces discussed so far---magnetism, the electric field, and gravity---all
arise from \emph{unbound} Aether: organized flow patterns, radial displacement, and
the residual asymmetry of transient binding.\ But the theory's central mechanism---the
ability of Aether to \emph{bind} into matter---has not yet been applied to the
strongest force in nature.

The strong nuclear force binds quarks into protons and neutrons, and protons and neutrons
into nuclei.\ Its defining features are well established experimentally: confinement (free
quarks are never observed), asymptotic freedom (quarks behave as free particles at short
distances), a force that is approximately constant with distance (the ``string tension''
$\sigma \approx 1$~GeV/fm), and spontaneous quark--antiquark pair creation when quarks
are pulled apart.\ In the Standard Model, these properties emerge from quantum
chromodynamics (QCD), a gauge theory with the symmetry group SU(3).

In the Aether framework, the strong force has a direct physical interpretation: it is the
elastic energy stored in a \emph{persistent bridge of compressed Aether} connecting quark
vortex endpoints.


\subsection{The Binding Bridge}

When Aether binds permanently (as described in Section~\ref{sec:core-binding}), the bound material no longer
flows or exerts pressure as a fluid---it is locked into mass.\ In the case of the electron,
binding produces a closed vortex ring whose stability rests on circulation quantization.\
But quark vortex endpoints can also anchor an open structure: a strand of highly compressed
Aether connecting two endpoints, held in its compressed state by the topological boundary
conditions the vortices impose.

In QCD, this structure is known as a \emph{color flux tube}---a narrow tube of
chromodynamic field energy connecting a quark to an antiquark (in a meson) or linking
three quarks at a Y-shaped junction (in a baryon).\ Lattice QCD calculations confirm
that the chromodynamic field is concentrated in such tubes rather than spreading
through all space.\ In the Aether model, the flux tube is not an abstract field
configuration but a physical object: a strand of compressed Aether with a definite
cross-sectional area, energy density, and material properties.

The energy of the bridge is proportional to its length:
%
\begin{equation}
  E_{\mathrm{bridge}} = \sigma \cdot L
  \label{eq:string-tension}
\end{equation}
%
where $\sigma$ is the string tension (energy per unit length) and $L$ is the bridge
length.\ This linear energy--distance relationship produces a \emph{constant} restoring
force:
%
\begin{equation}
  F_{\mathrm{strong}} = \frac{dE}{dL} = \sigma \approx 1\;\mathrm{GeV/fm}
\end{equation}
%
A force that does not weaken with distance is the signature of a material connection
rather than a field that spreads through space.\ The electric force weakens as $1/r^2$
because the radial flow dilutes over a sphere.\ The strong force does not weaken because
the bridge does not dilute---it is a one-dimensional material strand, not a
three-dimensional flow pattern.


\subsection{Confinement}

The confinement of quarks---the experimental fact that no free quark has ever been
observed---follows directly from the bridge model.\ When energy is supplied to separate
two quarks, the bridge stretches, storing energy $\sigma L$.\ When the stored energy
reaches the threshold for creating a new quark--antiquark pair---including both the
rest mass and the kinetic energy of the newly created quarks within the resulting
mesons, totalling $\sim$1~GeV---the bridge snaps.\ Each broken end
immediately forms a new vortex endpoint, producing two shorter bridges instead of
one long one.\ The result is two mesons (quark--antiquark pairs), not two free quarks.
Lattice QCD simulations confirm that this string-breaking transition occurs at a
quark separation of $r_{\mathrm{break}} \approx 1.2$~fm
\cite{bali2005,bulava2019}---consistent with $\sigma \times r_{\mathrm{break}}
\approx 1.2$~GeV after accounting for the short-distance Coulomb contribution
to the static potential.

This is exactly the mechanism of permanent Aether binding applied to a one-dimensional
geometry: when the energy density in the bridge exceeds the Dynamic Casimir Effect
threshold, new bound matter forms at the break point.\ The same DCE threshold that
governs electron creation governs the snap point of the bridge.

Confinement is therefore not a mysterious emergent property of a non-Abelian gauge
theory---it is the straightforward consequence of stretching a material strand until
it breaks and forms new endpoints.


\subsection{Asymptotic Freedom}

At very short distances, the strong coupling constant $\alpha_s$ decreases---quarks
inside a hadron interact more weakly when close together.\ This ``asymptotic freedom''
is one of the most counterintuitive features of QCD.

In the bridge model, asymptotic freedom has a geometric explanation rooted in the
bridge's compactness.\ The bridge diameter
$d_{\mathrm{bridge}} \approx 0.30$--$0.35$~fm is only $1$--$2$ times the healing
length $\xi \approx 0.17$--$0.28$~fm
(Section~\ref{sec:strong-constraints}).\ This means the bridge has no extended
``bulk interior''---the entire cross-section is within the healing-length
transition region, where the condensate density rises smoothly from the ambient
value to its peak and back.

When two quarks are separated by less than the bridge diameter
($L \lesssim d_{\mathrm{bridge}} \sim 1$--$2\,\xi$), the compressed
structure between them has not fully formed: the condensate perturbation is
still developing over the $\sim\!\xi$ recovery length on each side.\ The
quarks sit inside a partially-formed, low-amplitude perturbation rather than
a fully-compressed strand, so the effective restoring force is weak.\ As the
quarks separate and $L$ exceeds $d_{\mathrm{bridge}}$, the bridge becomes a
well-defined one-dimensional strand with full string tension $\sigma$, and
the confining force reaches its asymptotic constant value.

This predicts a specific crossover scale: the coupling should strengthen
rapidly as the quark separation passes through $\sim\!\xi \approx 0.2$~fm,
consistent with lattice QCD measurements showing the quark-antiquark potential
transitioning from approximately Coulombic ($V \propto 1/r$) at short range
to linear ($V = \sigma r$) at long range, with the crossover at
$\sim$0.2--0.5~fm.\ The crossover scale is not a free parameter---it is
determined by the healing length $\xi$, which is itself fixed by the
Aether quantum mass $m_A$.

Asymptotic freedom in this picture is not a running coupling constant in an
abstract renormalization group, but a geometric transition: from a region too
small for the condensate to form a coherent compressed structure (weak coupling)
to a fully-formed bridge (constant confining force).


\subsection{Gluon Self-Interaction}

In quantum electrodynamics, photons do not interact with each other---electromagnetic
waves in the Aether are linear (to the precision confirmed by experiment).\ In QCD,
gluons interact with each other, producing the non-Abelian structure that distinguishes
the strong force from electromagnetism.

In the bridge model, this distinction is natural.\ Photons are compression waves in
\emph{unbound} Aether, which is experimentally confirmed to be perfectly linear
(Hookean) in its pressure--density response.\ Gluons are excitations of the
\emph{compressed} bridge material---density waves, vibrational modes, and distortions
traveling along the compressed strand.\ Excitations in a highly compressed medium are
inherently nonlinear: phonons scatter off each other, solitons interact, and density
waves couple.\ The gluon self-interaction is not a postulate but the expected behavior
of disturbances propagating through a compressed material bridge.

This also explains why the strong force is non-Abelian while electromagnetism is
Abelian: the unbound Aether is a simple elastic medium (linear, Abelian), while the
compressed bridge interior is a dense, nonlinear system (non-Abelian).\ The
mathematical difference between the two forces reflects the physical difference
between wave propagation in a free medium and wave propagation in a material strand.


\subsection{The Origin of Hadron Mass}

The proton has a mass of 938~MeV/$c^2$, but the three valence quarks inside it
(two up quarks at $\sim$2.2~MeV and one down quark at $\sim$4.7~MeV) account for
only $\sim$9~MeV---less than 1\% of the total.\ The remaining 99\% is attributed in
QCD to ``gluon field energy'' and ``quark kinetic energy.''

In the bridge model, this mass has a concrete physical identity: it is the \emph{elastic
compression energy stored in the bridges}.\ The bridge material is Aether compressed
to $\sim\!7 \times 10^7$ times ambient density, and this compression energy is mass
($E = mc^2$).\ The proton's mass is predominantly the mass of the Y-shaped
bridge connecting its three quark endpoints, not the mass of the endpoints themselves.

This reinterpretation makes a quantitative prediction: the proton mass should be
calculable from the string tension $\sigma$, the bridge geometry (Y-shaped, with
three arms of length $\sim$0.3~fm meeting at a junction), and the endpoint (quark)
masses:
%
\begin{equation}
  m_p c^2 \approx 3 \times \sigma \times L_{\mathrm{arm}} + (2m_u + m_d)\,c^2 + E_{\mathrm{junction}}
\end{equation}
%
With $\sigma \approx 1$~GeV/fm and $L_{\mathrm{arm}} \approx 0.3$~fm, the bridge
contribution is $\sim$900~MeV, which is the correct order of magnitude.\ The
junction energy $E_{\mathrm{junction}}$ represents the binding energy at the point
where three bridges meet.


\subsection{Color as Bridge Orientation}

In QCD, quarks carry one of three ``color'' charges (red, green, blue), and observable
hadrons must be color-neutral.\ In the Aether framework, the three colors correspond
to the \emph{three independent spatial orientations} of the binding bridge.

A vortex line in three spatial dimensions has its circulation confined to a plane.
There are three independent planes (conventionally $xy$, $xz$, $yz$), and a binding
bridge inherits the orientation of the vortex circulation at its endpoints.\ The three
colors are these three orientations:
%
\begin{itemize}[leftmargin=2em]
  \item A bridge in the $xy$-plane carries ``red'' color.
  \item A bridge in the $xz$-plane carries ``green'' color.
  \item A bridge in the $yz$-plane carries ``blue'' color.
\end{itemize}
%
Color neutrality then has a geometric meaning: a baryon contains three bridges spanning
all three spatial planes, producing a composite structure that is rotationally
invariant---it has no preferred orientation.\ A meson contains a bridge and its
conjugate (a quark and an antiquark with opposite color orientations in the same plane),
which is also rotationally neutral in that plane.

The step from three orientations to SU(3) gauge symmetry requires the quantum
nature of the bridges.\ Classically, rotations among three planes give SO(3),
which has only 3~generators---far fewer than the 8~generators of SU(3).\ But
bridges in a quantum superfluid are quantum objects that can exist in coherent
superpositions of orientations:
%
\begin{equation}
  |\psi_{\mathrm{color}}\rangle = \alpha\,|xy\rangle + \beta\,|xz\rangle + \gamma\,|yz\rangle
\end{equation}
%
where $\alpha$, $\beta$, $\gamma$ are complex amplitudes.\ Unitary
transformations on this three-component complex vector space form U(3),
and after removing the overall phase, SU(3)---which has exactly
8~generators (the Gell-Mann matrices).\ In an isotropic medium all three
orientations are energetically equivalent, guaranteeing at minimum SO(3)
invariance.\ If the dynamics is also invariant under relative complex-phase
rotations between orientations---so that $|xy\rangle + |xz\rangle$ and
$|xy\rangle + i\,|xz\rangle$ have the same energy---the full SU(3)
symmetry is realized.

The eight gluons would then correspond to the eight independent ways of
unitarily mixing the three orientation amplitudes: six off-diagonal
generators (complex rotations between pairs of planes) and two diagonal
generators (relative phase differences).\ This geometric identification
parallels the proposal of Marsch and Narita \cite{marsch2021}, who
connect the three spin-helicity degrees of freedom in 3D to SU(3) color
structure.\ However, \emph{demonstrating that the nonlinear dynamics of
compressed Aether bridges respects the full SU(3) rather than only the
SO(3) subgroup remains an open problem}---it requires showing that all
eight collective excitation modes exist and that the bridge Lagrangian
is invariant under arbitrary unitary mixing of orientations.


\subsection{The Deconfinement Phase Transition}

At extreme temperatures ($T \gtrsim 2 \times 10^{12}$~K, corresponding to
$\sim$170~MeV), hadronic matter undergoes a phase transition to a quark--gluon plasma
(QGP), in which quarks and gluons move freely without confinement.\ This transition
has been observed at the Relativistic Heavy Ion Collider (RHIC) and the Large Hadron
Collider (LHC).

In the bridge model, the deconfinement transition is the \emph{thermal disruption of
the quark vortex endpoints that anchor the compressed bridges}.\ When thermal energy
exceeds the anchoring energy per unit length, the quark vortices destabilize and the
compressed strand relaxes back to ambient density.\ Without stable bridges,
there is no confinement, and quark vortex endpoints move freely through the medium.

The bridge material itself remains in the same thermodynamic phase as the surrounding
superfluid throughout this process (Section~\ref{sec:strong-constraints})---it is
simply compressed by the topological boundary conditions at the quark endpoints.\
The deconfinement temperature
$T_{\mathrm{deconf}} \approx 170$~MeV sets the energy scale of the binding:
%
\begin{equation}
  k_B T_{\mathrm{deconf}} \approx \sigma \cdot d_{\mathrm{bridge}}
\end{equation}
%
where $d_{\mathrm{bridge}}$ is the effective bridge diameter.\ With
$\sigma \approx 1$~GeV/fm, this gives $d_{\mathrm{bridge}} \approx 0.17$~fm---the
thermal energy scale of the bridge cross-section.\ This is smaller than the
geometric bridge diameter ($\sim$0.3--0.35~fm from the lattice flux tube profile,
Section~\ref{sec:strong-constraints}), because the deconfinement condition probes
the compact core of the bridge where the compression energy is concentrated, not
the full width including the healing-length tails ($\sim$0.4--0.6~fm).

The QGP state is therefore the Aether equivalent of a high-temperature plasma:
the compressed bridges have relaxed, and the formerly confined vortex endpoints (quarks)
move through a hot medium.\ As the plasma cools below $T_{\mathrm{deconf}}$,
quark vortices restabilize, bridges re-form, and quarks are re-confined.


\subsection{Quantum Fluctuations of the Bridge}
\label{sec:bridge-fluctuations}

The bridge model developed so far is classical: a static compressed strand
connecting quark vortex endpoints.\ A complete description requires
quantizing the fluctuations around this classical solution.\ The natural
framework is \emph{Bogoliubov--de~Gennes theory}---the same linearization
used in Section~\ref{sec:qss} for the homogeneous condensate, applied here
to the cylindrical bridge geometry.

The bridge supports three fluctuation families: \emph{longitudinal
phonons} (compression waves at $c_{\mathrm{phonon}} = c/\sqrt{X} \approx
10^{-4}\,c$, slow because $\Ka$ is constant while density is $X\rhoa$),
\emph{transverse oscillations} (string vibrations at speed $c$, since the
mass per unit length equals $\sigma/c^2$), and \emph{breathing modes}
(gapped radial oscillations, frozen out at low energy).\ The transverse
modes dominate the long-distance quantum correction.

\subsubsection*{The L\"uscher term}

The zero-point energy of the transverse fluctuation modes
produces a quantum correction to the quark--antiquark potential.\ For a
bridge of length $L$ with quarks at its endpoints (Dirichlet boundary
conditions), the allowed mode wavenumbers are
$k_n = n\pi / L$.\ The zero-point energy per transverse direction is
%
\begin{equation}
  E_{\mathrm{ZPE}} = \sum_{n=1}^{\infty} \frac{1}{2}\,\hbar\,\omega_n
    = \frac{\hbar\,c_s\,\pi}{2L} \sum_{n=1}^{\infty} n
\end{equation}
%
The divergent sum is regulated by zeta-function regularization,
$\sum n = \zeta(-1) = -1/12$, giving a finite contribution per
polarization:
%
\begin{equation}
  E_{\mathrm{ZPE}} = -\frac{\pi\,\hbar\,c_s}{24\,L}
\end{equation}
%
In three spatial dimensions the bridge has two transverse polarization
directions, so the total quantum correction to the static potential is
%
\begin{equation}
  \label{eq:luscher}
  \Delta V(L) = -\frac{\pi\,\hbar\,c_s}{12\,L}
\end{equation}
%
This is the bridge model's prediction of the \emph{L\"uscher term}
\cite{luscher1981,luscher2002}.\ In QCD, the L\"uscher term
$\Delta V = -\pi/(12\,r)$ (in natural units with $c = 1$) is a universal
quantum correction to the linear confining potential, confirmed by
high-precision lattice simulations \cite{juge2004}.\ The bridge model
reproduces the $-1/(12\,L)$ structure from the Bogoliubov
zero-point energy of a finite-length condensate strand.\ The coefficient
depends on the propagation speed of the transverse modes.\ The
\emph{longitudinal} (compression) modes propagate at
$c_{\mathrm{phonon}} = \sqrt{\Ka/\rho} = c/\sqrt{X} \approx 10^{-4}\,c$
inside the bridge, because the bulk modulus $\Ka$ is constant while the
density is $X\rhoa$.\ However, the L\"uscher term arises from
\emph{transverse} fluctuations, which in the Aether model are spin-wave
excitations of the condensate---the same mode that constitutes light
(Section~\ref{sec:qss}).\ The spin stiffness $K_{\mathrm{spin}}$
governing these transverse modes scales with density ($K_{\mathrm{spin}}
\propto \rho$, reflecting the stronger exchange coupling at closer
particle spacing), giving a spin-wave speed
$c_{\mathrm{spin}} = \sqrt{K_{\mathrm{spin}}/\rho} = c$ independent of
compression.\ This is the standard condensed matter distinction between
compression and spin-wave elastic moduli: the bulk modulus $\Ka$ (constant)
governs longitudinal modes, while the spin stiffness $K_{\mathrm{spin}}
\propto \rho$ governs transverse modes.\ The L\"uscher coefficient is
therefore set by $c_{\mathrm{spin}} = c$, matching the lattice-measured
value exactly.

\subsubsection*{String width broadening}

The transverse fluctuations cause the effective string width to
grow logarithmically with quark separation:
%
\begin{equation}
  \label{eq:string-width}
  w^2(L) = w_0^2 + \frac{\hbar\,c}{\pi\sigma}\,\ln\!\left(\frac{L}{L_0}\right)
\end{equation}
%
where $w_0 \sim \xi$ is the intrinsic bridge width, $L_0$ is a
short-distance cutoff of order $\xi$, and the coefficient $\hbar c/(\pi\sigma)$
sums over both transverse directions in $3{+}1$ dimensions.\ This
logarithmic broadening is a well-known prediction of the effective string
theory of confinement \cite{luscher2002}, and lattice simulations confirm
$w^2 \propto \ln L$ at large separations \cite{cardoso2013,cea2024}.\ At short separations
($L \lesssim 0.6$~fm), the measured width increases above the
asymptotic trend \cite{cea2024}, consistent with the bridge model's prediction
that asymptotic freedom dissolves the compressed structure when
$L \lesssim d_{\mathrm{bridge}}$.

\subsubsection*{Virtual pair creation and the pion cloud}

The most physically consequential fluctuation is \emph{virtual bridge
snapping}: a quantum fluctuation large enough to momentarily create a
quark--antiquark pair along the bridge, producing a short-lived meson
that is reabsorbed.\ The rate of such events is governed by the Schwinger
mechanism \cite{schwinger1951} applied to the bridge's chromoelectric
field.\ For a quark of mass $m_q$ in a field of string tension $\sigma$,
the pair-creation rate per unit length is
%
\begin{equation}
  \label{eq:schwinger-bridge}
  \frac{\Gamma}{L} \propto \sigma \,
    \exp\!\left(-\frac{\pi\,m_q^2\,c^3}{\hbar\,\sigma}\right)
\end{equation}
%
The exponent evaluates to $\pi m_q^2 c^3 / (\hbar\sigma) \approx
\pi m_q^2 / (\sigma \cdot \hbar c)$.\ Using $\hbar c \approx 197$~MeV\,fm
and $\sigma \approx 1$~GeV/fm $= 1000$~MeV/fm:
%
\begin{center}
\begin{tabular}{@{}lccl@{}}
\hline
\textbf{Quark} & $m_q$ (MeV) & $\pi m_q^2/(\sigma\hbar c)$ & \textbf{Suppression} \\
\hline
$u, d$ & 3--5   & $\sim 10^{-4}$ & None (copious) \\
$s$    & $\sim$95  & $\sim$0.14    & Mild \\
$c$    & $\sim$1275 & $\sim$26     & Strong \\
$b$    & $\sim$4180 & $\sim$280    & Exponential \\
\hline
\end{tabular}
\end{center}
%
For light quarks ($u$, $d$), the exponential is essentially unity---pair
creation is \emph{unsuppressed}.\ The bridge is constantly dressed by
a cloud of virtual pions (light $q\bar{q}$ bridge segments that briefly
form and recombine).\ For charm and heavier quarks, pair creation is
exponentially suppressed, and the bridge is a clean, well-defined string.

The crossover mass where pair creation turns off (Schwinger exponent
$\sim 1$) is
$m_q^* \sim \sqrt{\sigma\,\hbar c/\pi} \approx \sqrt{63{,}000}~\text{MeV}
\approx 250$~MeV---between the strange and charm quark masses.\ This
reproduces a fundamental feature of QCD: the transition from
light-quark (chiral) physics, dominated by pion-cloud effects, to
heavy-quark (potential model) physics, where the static string picture
is accurate.\ In the bridge model, this transition is not a change in
theoretical framework but a single physical parameter---the Schwinger
exponent---passing through unity.

\subsubsection*{Implications for the nucleon sigma term}

The pion cloud has a direct physical consequence: the nucleon sigma term
$\sigma_{\pi N} = \hat{m}\,\langle N|\bar{u}u + \bar{d}d|N\rangle$,
where $\hat{m} = (m_u + m_d)/2$.\ This measures how much of the
nucleon mass depends on the light quark masses.\ In the classical bridge
model (Section~\ref{sec:strong-force}), the quarks are ultra-relativistic:
their momentum ($p \sim \hbar/L_{\mathrm{arm}} \approx 660$~MeV/$c$)
vastly exceeds their rest mass ($m_q c \approx 5$~MeV/$c$), so the
bridge equilibrium geometry is insensitive to quark mass changes.\
The Hellmann--Feynman theorem gives a classical sigma term of only
$\sigma_{\pi N}^{\mathrm{classical}} \approx 3\hat{m}^2 c^2 / (pc)
\sim 0.06$~MeV---essentially zero.

The measured value, $\sigma_{\pi N} \approx 45$--$60$~MeV
\cite{alarcon2012,hoferichter2015}, is overwhelmingly dominated by
virtual pion exchange: the bridge fluctuations create and absorb
virtual pions whose mass depends on $m_q$ through the Gell-Mann--Oakes--Renner
relation ($m_\pi^2 \propto \hat{m}$).\ The bridge model thus predicts
the correct hierarchy: a negligible \emph{classical} sigma term (from
static geometry) plus a large \emph{quantum} contribution (from the
pion cloud dressed by Schwinger pair creation), with the quantum piece
dominating by two orders of magnitude.

Computing the full $\sigma_{\pi N}$ from the bridge model requires
calculating the pion--nucleon coupling constant $g_{\pi N}$ from the
overlap of the virtual pion's bridge wavefunction with the Y-junction
geometry of the nucleon---an open calculation that would provide a
stringent quantitative test.


\subsection{The Weak Force: Vortex Topology Transitions}

The weak nuclear force is responsible for processes that change one type of quark or
lepton into another---most notably beta decay, where a neutron transforms into a proton,
an electron, and an antineutrino.\ In the Standard Model, the weak force is mediated by
the massive $W^{\pm}$ and $Z^0$ bosons and is described by the SU(2) gauge symmetry.

In the Aether framework, the weak force is fundamentally different from the other three:
it is not a force in the traditional sense but a \emph{topology-changing transition}
of the vortex structure.\ When a down quark transforms into an up quark (emitting a
$W^-$ boson), the vortex endpoint changes its circulation configuration---the bridge
attached to it restructures.

The SU(2) gauge symmetry of the weak force arises from the spinor order parameter
$\Psi = (\psi_\uparrow, \psi_\downarrow)$ that is already required by the theory
for spin-statistics.\ The $W$ and $Z$ bosons are massive because they correspond to
\emph{gapped} excitations of the spinor condensate: changing the relative amplitude
or phase of the two spinor components costs energy proportional to the condensate
stiffness.\ Their large masses ($m_W \approx 80.4$~GeV, $m_Z \approx 91.2$~GeV)
reflect the enormous energy required to restructure the condensate at the scale
of a single vortex.

The short range of the weak force ($\sim 10^{-18}$~m) follows from the large boson
masses through the uncertainty principle: a virtual $W$ or $Z$ can propagate only
a distance $\sim \hbar / (m_W c) \approx 10^{-18}$~m before being reabsorbed.\ In
Aether terms, the topology-changing disturbance is heavily damped because it
requires restructuring the dense condensate, and it cannot propagate far before
the condensate reasserts its ground-state configuration.

Parity violation---the weak force's coupling only to left-handed particles---may
reflect a fundamental chirality in the vortex topology.\ If the binding bridge
attaches preferentially to one circulation handedness of the vortex endpoint,
the topology-changing transitions mediated by $W$ bosons would naturally couple
to only one chirality.\ This remains an open question requiring detailed analysis
of vortex--bridge junction geometry.


\subsection{The Four Forces Unified Through Aether Mechanics}

The four fundamental forces now have a unified interpretation within the Aether
framework:
%
\begin{center}
\begin{tabular}{@{}llll@{}}
\hline
\textbf{Force} & \textbf{Mechanism} & \textbf{Range} & \textbf{Strength} \\
\hline
Strong    & Compressed Aether bridge  & $\sim$1~fm (bridge snaps)    & $\sigma \approx 1$~GeV/fm \\
EM        & Flow vector of unbound Aether   & Infinite ($1/r^2$)    & $\alpha \approx 1/137$ \\
Weak      & Vortex topology transition      & $\sim 10^{-18}$~m     & $G_F \sim 10^{-5}$~GeV$^{-2}$ \\
Gravity   & Bind--unbind cycle asymmetry    & Infinite ($1/r^2$)    & $\alpha_G \sim 10^{-39}$ \\
\hline
\end{tabular}
\end{center}
%
The force hierarchy is natural:
%
\begin{itemize}[leftmargin=2em]
  \item The strong force is the strongest because it involves the \emph{full binding
    energy} of condensed Aether---the most energetic process available in the medium.
  \item The electromagnetic force is weaker because it involves the \emph{pressure of
    flowing} unbound Aether---significant, but not binding-scale.
  \item The weak force is weaker still because it requires \emph{restructuring the
    condensate}---a costly process damped by the large condensate stiffness.
  \item Gravity is the weakest because it is only the \emph{tiny residual asymmetry}
    of the transient bind--unbind cycle---the faintest whisper of the binding process.
\end{itemize}
%
All four forces arise from the same medium and its mechanical capabilities.\
The strong force is elastic compression anchored by quark vortex endpoints.\
Electromagnetism is the flow pattern of unbound Aether.\ The weak force is the
topological restructuring of vortex endpoints.\ Gravity is the leftover asymmetry
of transient binding.\ There is no need for separate, unrelated mechanisms---the
four forces are four aspects of one medium.


% ============================================================

% ============================================================
\section{Summary of Theoretical Properties}
% ============================================================

\subsection{Properties of the Strong Force}

\begin{enumerate}[leftmargin=2em]
  \item Quark confinement arises from compressed Aether bridges---condensate
    strands connecting quark vortex endpoints, with string tension
    $\sigma \approx 1$~GeV/fm set by the elastic compression energy of the
    logarithmic equation of state.
  \item Bridges are elastic compressions of the same superfluid phase (compression
    ratio $X \approx 7 \times 10^7$), not a separate bound phase.\ Deconfinement
    at $T \approx 170$~MeV corresponds to thermal disruption of the anchoring quark
    vortices, not melting of a distinct phase.
  \item Asymptotic freedom follows from bridge compactness: $d_{\mathrm{bridge}}/\xi
    \approx 1$--$2$ means bridges are ``all core'' with no bulk interior.\ At
    separations $L < d_{\mathrm{bridge}}$, the condensate perturbation has not fully
    formed, giving weak coupling; at $L > d_{\mathrm{bridge}}$, full string tension
    emerges.\ The Coulomb-to-linear crossover at $L \sim \xi \approx 0.2$~fm matches
    lattice QCD with no free parameters.
  \item Color charge corresponds to quantum superpositions of three bridge
    orientations (xy, xz, yz planes); unitary mixing of these three complex
    amplitudes gives SU(3) symmetry with 8~generators.\ Full SU(3) invariance
    (beyond the SO(3) subgroup) requires the bridge dynamics to be invariant
    under relative complex-phase rotations between orientations.
  \item Hadron mass is predominantly bridge energy ($\sigma \times L$), not quark
    rest mass: three bridge arms of length $\sim$0.3~fm give $3\sigma L_{\mathrm{arm}}
    \approx 900$~MeV, reproducing the proton mass with quark and junction corrections.
  \item Quantum fluctuations of the bridge produce three mode families: longitudinal
    phonons ($c_{\mathrm{phonon}} = c/\sqrt{X}$ at compression $X$, from the
    constant-$\Ka$ logarithmic EOS), transverse spin-wave oscillations
    ($c_{\mathrm{spin}} = c$, from $K_{\mathrm{spin}} \propto \rho$),
    and gapped breathing modes.
  \item The L\"uscher term $\Delta V = -\pi\hbar c_s/(12L)$ emerges from the
    zero-point energy of two transverse polarizations, reproducing the universal
    $-1/(12\,r)$ quantum correction to the confining potential.
  \item Schwinger pair creation along the bridge produces a pion cloud for light
    quarks (exponent $\ll 1$) but is exponentially suppressed for heavy quarks.\ The
    crossover mass $m_q^* \approx 250$~MeV naturally separates chiral QCD from
    potential-model QCD.
  \item String-breaking at $L \approx 1.2$~fm (matching lattice QCD) occurs when
    the bridge stores enough energy for quark--antiquark pair creation including
    kinetic energy.
  \item The flux tube width ($w \approx 0.5$~fm) matches the bridge model prediction
    $2$--$3\,\xi \approx 0.4$--$0.6$~fm with no adjustable parameters.
  \item The weak force arises from vortex topology transitions (core reconnection
    of quark vortex endpoints), with the short range ($\sim 10^{-18}$~m) set by the
    massive $W$ and $Z$ bosons as gapped excitations of the spinor condensate.
\end{enumerate}




% ============================================================
\section{Experimental Support and Connections to Established Physics}
% ============================================================

The following experimentally confirmed phenomena support or align with the aspects of the theory presented in this paper:

\begin{enumerate}[leftmargin=2em]
  \item \textbf{Lattice QCD string-breaking distance.}\ Lattice simulations
    with dynamical quarks \cite{bali2005,bulava2019} observe the static
    quark--antiquark potential levelling off at a separation
    $r_{\mathrm{break}} \approx 1.2$~fm, confirming string breaking by
    light-quark pair creation.\ The bridge model predicts this distance
    from the string tension $\sigma \approx 1$~GeV/fm and the
    pair-creation energy threshold ($\sim$1~GeV, including rest mass and
    kinetic energy of the created quarks), giving
    $L_{\mathrm{snap}} \approx 1.2$~fm---in direct quantitative agreement
    with the lattice measurement.

  \item \textbf{Lattice QCD flux tube width.}\ High-precision lattice
    simulations \cite{cardoso2013,cea2024} measure the transverse
    chromoelectric field width at $w \approx 0.5$~fm for quark separations
    $d \gtrsim 0.7$~fm.\ The bridge model predicts the compressed region
    spans $\sim\!2$--$3\,\xi \approx 0.4$--$0.6$~fm, in quantitative
    agreement with no adjustable parameters.\ Earlier dual-superconductor
    fits \cite{cea2012} extract a penetration depth
    $\lambda \approx 0.22$--$0.24$~fm that overlaps with the Aether healing
    length $\xi \approx 0.17$--$0.28$~fm.

  \item \textbf{L\"uscher term in the static quark potential.}\ Lattice QCD
    simulations \cite{juge2004,luscher2002} confirm that the quark--antiquark
    potential at large separations deviates from pure linearity by a
    universal quantum correction $\Delta V = -\pi/(12\,r)$.\ The bridge
    model derives this term from the zero-point energy of transverse
    Bogoliubov modes confined to a finite-length condensate strand
    (Section~\ref{sec:bridge-fluctuations}), reproducing the $-1/(12\,r)$
    structure with the coefficient set by the transverse (spin-wave) speed
    $c_{\mathrm{spin}} = c$, which is density-independent because the spin
    stiffness scales linearly with density ($K_{\mathrm{spin}} \propto \rho$).

  \item \textbf{Light-quark versus heavy-quark phenomenology.}\ QCD exhibits
    a sharp transition: light quarks ($u$, $d$, $s$) are governed by chiral
    dynamics and strong pion-cloud effects, while heavy quarks ($c$, $b$)
    are accurately described by static potential models.\ The bridge model
    reproduces this transition from a single mechanism: the Schwinger
    pair-creation exponent $\pi m_q^2/(\sigma\hbar c)$ passes through
    unity at $m_q^* \approx 250$~MeV, between the strange and charm masses,
    naturally separating copious pion-cloud dressing (light quarks) from
    clean string dynamics (heavy quarks).

  \item \textbf{Asymptotic freedom.}\ Deep inelastic scattering experiments
    confirm that quarks interact more weakly at short distances.\ In the bridge
    model this is geometric: when quarks are close, the compressed Aether bridge
    between them is short and nearly all core ($d/\xi \approx 1$--$2$), so there
    is little compressed material to resist further approach.

  \item \textbf{Hadron mass from bridge energy.}\ The proton mass
    ($\approx 938$~MeV) far exceeds the summed current quark masses
    ($\approx 10$~MeV).\ The bridge model attributes the excess to
    compressive energy stored in the Aether bridges: string tension times bridge
    length ($\sigma \times L \approx 900$~MeV) reproduces the measured mass with
    quark and junction corrections.

  \item \textbf{The deconfinement phase transition.}\ Heavy-ion collisions at
    RHIC and the LHC produce quark--gluon plasma at $T \approx 170$~MeV, in
    which quarks move freely.\ In the bridge model, this is thermal disruption
    of compressed bridges: when thermal energy exceeds the bridge binding
    energy, the compressed Aether channels dissolve and quarks deconfine.

\end{enumerate}


% ============================================================
\section{Open Questions and Testable Predictions}
% ============================================================

The following areas relevant to this paper remain underdeveloped and represent opportunities for further refinement or falsification:

\begin{enumerate}[leftmargin=2em]
  \item \textbf{Nucleon sigma term from bridge fluctuations.}\ The classical
    bridge model predicts a negligible sigma term
    ($\sigma_{\pi N}^{\mathrm{classical}} \sim 0.1$~MeV) because the quarks
    are ultra-relativistic and the bridge geometry is insensitive to their
    mass.\ The measured $\sigma_{\pi N} \approx 45$--$60$~MeV
    \cite{alarcon2012,hoferichter2015} is dominated by the pion cloud
    arising from Schwinger pair creation along the bridge
    (Section~\ref{sec:bridge-fluctuations}).\ Computing the full sigma term
    requires deriving the pion--nucleon coupling $g_{\pi N}$ from the
    overlap of the virtual-pion bridge wavefunction with the Y-junction
    geometry---a calculation that would test the quantum bridge framework
    against one of the most precisely known hadronic quantities.

\end{enumerate}



% ============================================================
\section{Mathematical Summary}
% ============================================================

This section collects the key equations relevant to the topics covered in this paper.
\noindent\textit{For the complete mathematical summary, see Paper~1: General Overview \cite{otto-overview}.}

\subsection{Fundamental Aether Parameters}
% --------------------------------------------------

The theory posits four microscopic parameters from which all macroscopic
phenomena derive:
%
\begin{center}
\begin{tabular}{@{}llll@{}}
\hline
\textbf{Symbol} & \textbf{Name} & \textbf{Constrained range} & \textbf{Primary constraint} \\
\hline
$\rhoa$  & Ambient density     & $10^{11}$--$10^{12}$~kg/m$^3$         & Electron vortex energy \\
$\Ka$    & Bulk modulus         & $10^{28}$--$10^{29}$~J/m$^3$          & $\Ka = \rhoa\,c^2$ \\
$m_A$    & Aether quantum mass  & $500$--$850$~MeV/$c^2$                & Flux tube width / glueball mass \\
$\xi$    & Healing length       & $0.17$--$0.28$~fm                     & Bridge diameter (lattice QCD) \\
\hline
\end{tabular}
\end{center}

These four are not independent.\ Two defining relations connect them:
%
\begin{align}
  c &= \sqrt{\Ka / \rhoa}
    \label{eq:sum-speed-of-light} \\
  \xi &= \frac{\hbar}{m_A\,c\,\sqrt{2}}
    \label{eq:sum-healing-length}
\end{align}
%
leaving two free Aether parameters (e.g., $\rhoa$ and $m_A$) from which all
others follow.\ (A third parameter, $\epsilon$, enters through gravity; see below.)


% --------------------------------------------------
\subsection{Equation of State}
% --------------------------------------------------

The Aether obeys a logarithmic equation of state, the unique form that produces
a density-independent bulk modulus (Section~\ref{sec:qss}):
%
\begin{equation}
  P = \Ka \ln\!\Bigl(\frac{\rho}{\rho_0}\Bigr) + P_0
  \label{eq:sum-eos}
\end{equation}
%
Key consequence: the speed of sound $c_s = \sqrt{\Ka/\rho}$ varies with
density.\ At ambient density, $c_s = c$.\ Note: in gravitational fields,
Bernoulli's equation for free-fall inflow gives $\Delta\rho/\rhoa = 0$
exactly, so $c$ does not vary near gravitating masses; the density
dependence is relevant for material media and for the two-fluid
cosmological structure.\ This equation of state guarantees that $\Ka$ is
constant under compression, which is experimentally confirmed by:
%
\begin{itemize}[leftmargin=2em]
  \item Velocity time dilation measurements spanning $10^{-6}\,c$ to $0.9994\,c$
    (precision $\sim 10^{-8}$).
  \item The Fresnel drag coefficient $1 - 1/n^2$, exact across materials with
    density ratios up to $\sim$6:1.
\end{itemize}


% --------------------------------------------------
\subsection{Strong Force (Bound Aether Bridges)}
% --------------------------------------------------

The strong force is a physical strand of compressed Aether connecting
quark vortex endpoints (Section~\ref{sec:strong-force}).
%
\paragraph{String tension.}
%
\begin{equation}
  \sigma \approx 1\;\mathrm{GeV/fm} = 1.6 \times 10^5\;\mathrm{J/m}
  \label{eq:sum-string-tension}
\end{equation}

\paragraph{Bridge diameter from deconfinement.}
%
\begin{equation}
  d_{\mathrm{bridge}} = \frac{k_B\,T_{\mathrm{deconf}}}{\sigma}
    \approx 0.17\;\mathrm{fm}
  \label{eq:sum-bridge-diameter}
\end{equation}

\paragraph{Bridge energy density.}
%
\begin{equation}
  \varepsilon_{\mathrm{bridge}} = \frac{\sigma}{A}
    = \frac{4\sigma}{\pi\,\xi^2}
    \sim 10^{36}\text{--}10^{37}\;\mathrm{J/m^3}
  \label{eq:sum-bridge-eps}
\end{equation}

\paragraph{Compression ratio.}
The ratio of bridge to ambient density is determined by the three
constrained quantities $\sigma$, $\xi$, and $\Ka$
(Section~\ref{sec:strong-constraints}):
%
\begin{equation}
  X \equiv \frac{\rho_{\mathrm{bridge}}}{\rhoa}
    = \frac{4\sigma}{\pi\,\xi^2\,\Ka}
    \approx 7 \times 10^7
  \label{eq:sum-compression}
\end{equation}
%
The compression energy from the logarithmic equation of state:
%
\begin{equation}
  \varepsilon_{\mathrm{compress}} = \Ka\bigl[X - \ln X - 1\bigr]
    \approx \Ka\,X \quad (\text{for } X \gg 1)
  \label{eq:sum-compress-energy}
\end{equation}

\paragraph{Proton mass from bridge geometry.}
Three Y-junction arms of length $\sim$0.3~fm:
%
\begin{equation}
  m_p\,c^2 \approx 3\,\sigma\,L_{\mathrm{arm}} + m_{\mathrm{quarks}}\,c^2
    + E_{\mathrm{junction}}
    \approx 900 + 9 + 29 = 938\;\mathrm{MeV}
  \label{eq:sum-proton}
\end{equation}
%
where $E_{\mathrm{junction}} \approx 29$~MeV is the residual binding
energy at the Y-shaped junction, inferred from the difference between the
measured proton mass and the bridge plus quark contributions.

\paragraph{Wave speeds in bridge.}
The bridge has two distinct propagation speeds:
%
\begin{align}
  c_{\mathrm{phonon}} &= \frac{c}{\sqrt{X}}
    \approx 10^{-4}\,c
    \approx 3 \times 10^4\;\mathrm{m/s}
    \quad\text{(longitudinal, from constant $\Ka$)}
    \label{eq:sum-bridge-cphonon} \\
  c_{\mathrm{spin}} &= c
    \quad\text{(transverse, from $K_{\mathrm{spin}} \propto \rho$)}
    \label{eq:sum-bridge-cspin}
\end{align}
%
The L\"uscher term coefficient is set by $c_{\mathrm{spin}} = c$,
matching the lattice QCD value $\Delta V = -\pi/(12\,r)$.


% --------------------------------------------------

\section*{Acknowledgments}
\addcontentsline{toc}{section}{Acknowledgments}

The author acknowledges the use of Claude (Anthropic) for \LaTeX{} formatting
and editorial assistance.

\bigskip


% ============================================================
% BIBLIOGRAPHY
% ============================================================

\begin{thebibliography}{99}

\bibitem{schwinger1951}
Schwinger, J.,
``On Gauge Invariance and Vacuum Polarization,''
\textit{Physical Review}, \textbf{82}, 664--679 (1951).


\bibitem{cardoso2013}
Cardoso, N., Cardoso, M., and Bicudo, P.,
``Inside the SU(3) quark-antiquark QCD flux tube: screening versus quantum widening,''
\textit{Phys.\ Rev.\ D}, \textbf{88}(5), 054504 (2013).


\bibitem{marsch2021}
Marsch, E.\ and Narita, Y.,
``Threefold spin helicity as possible origin of SU(3) gauge symmetry,''
\textit{Eur.\ Phys.\ J.\ Plus}, \textbf{136}, 652 (2021).


\bibitem{bali2005}
Bali, G.~S., Neff, H., Duessel, T., Lippert, T., and Schilling, K.,
``Observation of string breaking in QCD,''
\textit{Phys.\ Rev.\ D}, \textbf{71}(11), 114513 (2005).


\bibitem{bulava2019}
Bulava, J., H\"orz, B., Knechtli, F., Koch, V., Moir, G., Morningstar, C., and Peardon, M.,
``String breaking by light and strange quarks in QCD,''
\textit{Phys.\ Lett.\ B}, \textbf{793}, 493--498 (2019).


\bibitem{cea2024}
Cea, P., Cosmai, L., and Ferretti, A.,
``The chromoelectric flux tube revisited,''
arXiv:2409.20168 [hep-lat] (2024).


\bibitem{cea2012}
Cea, P.\ and Cosmai, L.,
``Chromoelectric flux tubes and coherence length in QCD,''
\textit{Phys.\ Rev.\ D}, \textbf{86}(5), 054501 (2012).



\bibitem{luscher1981}
L\"uscher, M.,
``Symmetry-breaking aspects of the roughening transition in gauge theories,''
\textit{Nuclear Physics B} \textbf{180}, 317--329 (1981).


\bibitem{luscher2002}
L\"uscher, M. and Weisz, P.,
``Quark confinement and the bosonic string,''
\textit{Journal of High Energy Physics} \textbf{2002}(07), 049 (2002).


\bibitem{juge2004}
Juge, K.~J., Kuti, J., and Morningstar, C.,
``Fine structure of the QCD string spectrum,''
\textit{Physical Review Letters} \textbf{90}, 161601 (2003).


\bibitem{alarcon2012}
Alarcón, J.~M., Martin Camalich, J., and Oller, J.~A.,
``The chiral representation of the $\pi N$ scattering amplitude and the pion-nucleon sigma term,''
\textit{Physical Review D} \textbf{85}, 051503(R) (2012).


\bibitem{hoferichter2015}
Hoferichter, M., Ruiz de Elvira, J., Kubis, B., and Meißner, U.-G.,
``High-precision determination of the pion-nucleon $\sigma$ term from Roy--Steiner equations,''
\textit{Physical Review Letters} \textbf{115}, 092301 (2015).




\bibitem{otto-overview}
Otto, E.,
``Unifying Theory of Dynamic Aether Mechanics: General Overview,''
(2026).
\end{thebibliography}

\end{document}